% Standalone problem document, generated from the adjacent README.md.
% Compile with XeLaTeX. Mathematical content is maintained in Markdown.
\documentclass[11pt,a4paper]{article}
\usepackage[fontsize=10.5pt]{scrextend}
\usepackage[margin=22mm,headheight=15pt,headsep=9mm,footskip=11mm]{geometry}
\usepackage{fontspec}
\setmainfont{texgyrepagella}[Extension=.otf,UprightFont=*-regular,BoldFont=*-bold,ItalicFont=*-italic,BoldItalicFont=*-bolditalic]
\setsansfont{texgyreheros}[Extension=.otf,UprightFont=*-regular,BoldFont=*-bold,ItalicFont=*-italic,BoldItalicFont=*-bolditalic]
\setmonofont{texgyrecursor}[Extension=.otf,UprightFont=*-regular,BoldFont=*-bold,ItalicFont=*-italic,BoldItalicFont=*-bolditalic,Scale=MatchLowercase]
\usepackage{amsmath,amssymb}
\usepackage{unicode-math}
\setmathfont{texgyrepagella-math.otf}
\usepackage{xcolor,microtype,enumitem,needspace,fancyhdr}
\usepackage{bookmark}
\definecolor{ink}{HTML}{17344B}
\definecolor{accent}{HTML}{197D80}
\definecolor{muted}{HTML}{596773}
\hypersetup{colorlinks=true,linkcolor=accent,urlcolor=accent,pdftitle={FR-10 - Sharp
sampling complexity for Walsh restricted
isometries},pdfauthor={Open Problems in Numerical Linear Algebra}}
\urlstyle{same}
\setlength{\parindent}{0pt}
\setlength{\parskip}{5pt plus 1pt minus 1pt}
\linespread{1.04}
\setlength{\emergencystretch}{3em}
\widowpenalty=10000
\clubpenalty=10000
\displaywidowpenalty=10000
\allowdisplaybreaks[1]
\setlist{nosep,leftmargin=1.5em}
\providecommand{\tightlist}{\setlength{\itemsep}{0pt}\setlength{\parskip}{0pt}}
\providecommand{\pandocbounded}[1]{#1}
\pagestyle{fancy}
\fancyhf{}
\fancyhead[L]{\sffamily\footnotesize\color{muted}OPEN PROBLEMS IN NUMERICAL LINEAR ALGEBRA}
\fancyhead[R]{\sffamily\bfseries\footnotesize\color{accent}FR-10}
\fancyfoot[L]{\parbox[b]{0.9\textwidth}{\sffamily\fontsize{8}{10}\selectfont\color{muted}Editorial ratings; a bounded search cannot certify that a problem remains unsolved.\\Literature check: 2026-09-08}}
\fancyfoot[R]{\sffamily\footnotesize\color{muted}\thepage}
\renewcommand{\headrulewidth}{0.3pt}
\renewcommand{\headrule}{\hbox to\headwidth{\color{accent}\leaders\hrule height \headrulewidth\hfill}}
\makeatletter
\renewcommand{\section}{\Needspace{5\baselineskip}\@startsection{section}{1}{0pt}{12pt plus 2pt minus 2pt}{4pt}{\sffamily\bfseries\large\color{ink}}}
\renewcommand{\subsection}{\Needspace{4\baselineskip}\@startsection{subsection}{2}{0pt}{9pt plus 1pt minus 1pt}{3pt}{\sffamily\bfseries\normalsize\color{ink}}}
\makeatother
\setcounter{secnumdepth}{0}
\begin{document}
{\sffamily\small\color{muted}Frames and matrix designs\par}
\vspace{3pt}
{\raggedright\hyphenpenalty=10000\exhyphenpenalty=10000\sffamily\bfseries\fontsize{21}{24}\selectfont\color{ink}Sharp
sampling complexity for Walsh restricted isometries\par}
\vspace{7pt}
{\sffamily\small\textbf{Difficulty:} challenging\quad\textbf{Importance:} broadly
interesting\par}
\vspace{4pt}
{\color{accent}\hrule height 0.8pt}
\vspace{6pt}
\textbf{Status:} source-backed open problem; no resolution found in the
bounded literature check.

Let \(N=2^d\) with \(d\ge1\) and index rows and columns by
\(\mathbb F_2^d\). The normalized Walsh matrix is \[
(H_N)_{a,b}=N^{-1/2}(-1)^{a\cdot b}.
\] For \(m\ge1\), choose row indices \(a_1,\ldots,a_m\) independently
and uniformly with replacement and set
\(\Phi=\sqrt{N/m}(H_N)_{(a_1,\ldots,a_m),:}\). Let \(E_{m,N,k}\) be the
event that \[
\tfrac12\|x\|_2^2\le\|\Phi x\|_2^2\le\tfrac32\|x\|_2^2
\quad\text{for every }x\in\mathbb R^N\text{ with }|\operatorname{supp}x|\le k.
\] Define \[
m_*(N,k)=\min\{m\in\mathbb N:m\ge1,\ \Pr(E_{m,N,k})\ge0.9\}.
\] \textbf{Problem.} Determine \(m_*(N,k)\) up to universal
multiplicative constants, uniformly for \(1\le k\le N\), including all
necessary logarithmic factors. The same sample must work for all sparse
vectors. This is a quantitative restatement of the published sampling
gap, not a separately numbered conjecture; the success probability and
distortion fix a convention.

The Walsh transform is the Fourier transform on \(\mathbb F_2^d\). Its
subspace structure differs from the cyclic Fourier ensemble in FR-02, so
their lower bounds cannot be interchanged.

\subsection{References}\label{references}

\begin{enumerate}
\def\labelenumi{\arabic{enumi}.}
\tightlist
\item
  J. Błasiok, P. Lopatto, K. Luh, J. Marcinek, and S. Rao, \emph{An
  improved lower bound for sparse reconstruction from subsampled Walsh
  matrices}, Discrete Analysis 2023:3, introduction and Theorem 3.1.
  \href{https://arxiv.org/abs/1903.12135}{Paper}.
\item
  I. Haviv and O. Regev, \emph{The restricted isometry property of
  subsampled Fourier matrices}, Theorem 1.1 and its
  bounded-orthonormal-matrix scope.
  \href{https://arxiv.org/abs/1507.01768}{Paper}.
\end{enumerate}

\subsection{Status check --- 2026-09-08}\label{status-check-2026-09-08}

The lower-bound paper uses independent Bernoulli row inclusion; the
upper-bound literature also treats independent draws with replacement,
the explicit convention here. Its \(\Omega(k\log k\log(N/k))\)
obstruction applies in a specified intermediate sparsity range, not as
an all-parameter formula. The logarithmic gap discussed there remains
unclosed in the targeted search for later Walsh RIP results. Endpoint
regimes must also be accounted for; in particular \(m_*(N,1)=1\).
\end{document}
